% paper/skeleton.tex
% ───────────────────────────────────────────────────────────────────
% Kavach: A Parliament-of-Ministers Semantic Firewall for AI Agents
% Working title — fill in numbers as workstream D produces them.
%
% Target venues (in order of submission preference):
%   1. ICML 2026 Workshop on Agents in the Wild — non-archival, May/June
%   2. MASEC @ NeurIPS 2026 — non-archival workshop, September
%   3. SaTML 2027 — archival, September-October
%   4. DEF CON AI Village 2026 — talk, June
% ───────────────────────────────────────────────────────────────────

\documentclass[11pt]{article}

\usepackage[margin=1in]{geometry}
\usepackage{amsmath, amssymb, amsthm}
\usepackage{algorithm, algorithmic}
\usepackage{booktabs}
\usepackage{hyperref}
\usepackage{graphicx}
\usepackage{xcolor}
\usepackage{cite}

\newtheorem{theorem}{Theorem}
\newtheorem{definition}{Definition}
\newtheorem{claim}{Claim}

\title{Kavach: A Parliament-of-Ministers Semantic Firewall \\
       for Real-Time Interception of LLM Agent Tool Calls}
\author{[Author names — PES University]}
\date{\today}

\begin{document}
\maketitle

% ═══════════════════════════════════════════════════════════════════
\begin{abstract}
We present Kavach, a runtime guardrail for LLM agents that intercepts every
tool invocation before execution, evaluates it against a curated corpus of
attack-pattern descriptions across four semantic ministers
(EXECUTOR / VAULT / CHANNEL / NAVIGATOR), and combines minister verdicts via
a deterministic speaker into a BLOCK / ESCALATE / ALLOW decision in a
real-time budget of [TBD]\,ms p95. We additionally introduce COMPASS, a
session-level alignment oracle that flags drift between the user's stated
intent and the agent's proposed actions. Evaluated on InjecAgent
(1{,}054 cases) and a hand-curated benign trace set (50 sessions), Kavach
achieves [TBD] BLOCK-recall on attacks at [TBD] FPR on benign work,
representing [TBD]× reduction in attack success rate relative to
LlamaFirewall \cite{saxe2025llamafirewall} and AGrail
\cite{luo2025agrail} at comparable utility. We frame Kavach's primary
architectural contribution as the \emph{temporal--spatial} defense-in-depth
distinction: existing guardrails enforce in parallel layers at the same
time-step; Kavach enforces in serialized layers ordered by trust acquisition,
ensuring no agent extension code can run before the ministers have validated
it.
\end{abstract}

% ═══════════════════════════════════════════════════════════════════
\section{Introduction}
\label{sec:intro}

LLM agents have reached the operational sophistication that production
systems delegate non-trivial work to them — code generation, infrastructure
operations, customer support, sales pipelines. The agent's harness — the
deterministic code that wraps tool calls, manages context, and persists state
— now does the bulk of the engineering work; the model itself accounts for
roughly 1.6\% of source lines in production systems
\cite{liu2026clawdesign}. This shift has produced a new and growing class of
vulnerabilities: the harness, not the model, is the attack surface.

Recent CVEs make the pattern unambiguous. CVE-2025-59536 and
CVE-2026-21852 against Claude Code's pre-trust hooks; CVE-2025-68664 against
LangChain's serialization paths (98M monthly installs at disclosure time);
CVE-2025-34291 against Langflow's CORS handling. These are not model
failures. They are infrastructure failures, and they are exploited by
infrastructure-shaped attacks: prompt-injection in tool outputs, credential
exfiltration via metadata services, persistence install via lifecycle hooks.

We argue that the right defense for this surface is a runtime semantic
firewall — a deterministic layer that evaluates every tool invocation
against a curated body of attack-pattern semantics before the tool executes.
This paper presents Kavach, a working implementation of that idea.

\paragraph{Contributions.} This paper makes three contributions:
\begin{enumerate}
    \item A \emph{parliament-of-ministers} architecture for runtime agent
          tool-call evaluation. Four specialized ministers — EXECUTOR for
          code execution, VAULT for credential theft, CHANNEL for data
          exfiltration, NAVIGATOR for trajectory drift — query disjoint
          ChromaDB collections of attack-pattern descriptions and produce
          confidence-bounded verdicts that a deterministic speaker
          aggregates into a final BLOCK / ESCALATE / ALLOW outcome
          (\S\ref{sec:architecture}).

    \item A formalization of the \emph{temporal--spatial} defense-in-depth
          distinction. Existing guardrails — LlamaFirewall, AgentSpec, AGrail,
          ShieldAgent — enforce in spatial layers (parallel rails at one
          time-step). Kavach additionally enforces in temporal layers
          (ordered by trust acquisition: ministers must be online and
          authoritative before any agent extension code can register).
          We give a formal claim about the failure modes each layering
          discipline can and cannot cover (\S\ref{sec:temporal-spatial}).

    \item An evaluation methodology with explicit safeguards against
          training-set contamination. We ran the InjecAgent benchmark
          \emph{cold} — the corpus was authored from MITRE ATT\&CK technique
          IDs and OWASP Agentic 2026 categories, never from any specific CVE
          or test case. We measured benign FPR before measuring attack
          recall, treating FPR\,$<$\,5\% on hand-curated benign traces as a
          gate condition. We report cost (latency, embedding queries per
          decision) jointly with detection numbers, following Kapoor et
          al.~\cite{kapoor2024aiagents}.
\end{enumerate}

% ═══════════════════════════════════════════════════════════════════
\section{Background}
\label{sec:background}

\subsection{Agent infrastructure attack surface}
% Cover: tool calls, MCP, the harness shift, the 1.6/98.4 inversion
% Cite: Liu 2604.14228, Hou 2506.02040 ETDI 2506.01333 STRIDE 2603.22489

\subsection{Existing guardrails}
% Layered taxonomy of: prompt-time (PromptGuard 2), CoT-time (AlignmentCheck),
% tool-time (LlamaFirewall, AgentSpec, AGrail, ShieldAgent),
% reply-time (NeMo output rails, content safety APIs).
% Cite: 2505.03574, 2503.18666, 2502.11448, 2503.22738

\subsection{What's missing}
% Three gaps: (1) tool-time enforcement assumes the harness is itself trusted
% — Kavach addresses this in §4; (2) the speaker in existing parliaments is
% a fixed combiner rather than a generator-evaluator with a deterministic
% backstop; (3) no published guardrail reports cost-controlled evaluation.

% ═══════════════════════════════════════════════════════════════════
\section{Kavach Architecture}
\label{sec:architecture}

\subsection{Overview}
% Pipeline figure: user msg → seed_intent → action proposed →
%   COMPASS drift check → semantic router → activated ministers in parallel →
%   speaker → BLOCK/ESCALATE/ALLOW
% Cite our own openclaw.plugin.json wire format.

\subsection{Ministers}
% Per minister: source taxonomy, document count, embedding model, threshold.
% Table: minister × (#patterns, #docs, sources cited).

\subsubsection{EXECUTOR}
% MITRE T1059, T1190, T1195, T1546, T1574; OWASP A09.

\subsubsection{VAULT}
% MITRE T1552, T1078, T1539, T1555; OWASP A02, A06.

\subsubsection{CHANNEL}
% MITRE T1041, T1567, T1071, T1048, T1572; OWASP A04.

\subsubsection{NAVIGATOR}
% MITRE T1083, T1057, T1005, T1135, T1119; OWASP A01.

\subsection{COMPASS — alignment oracle}
% Stage II seed_intent stores BGE vector of user message.
% At each tool call, cosine(intent, action). Drift threshold from Youden's J.
% Maps to Saha et al.'s task-alignment property in their formal framework.
% Cite: 2603.19469.

\subsection{Speaker — deterministic verdict combiner}
% Five cases as in speaker.py: blocking minister wins; drift+escalate→block;
% any escalate→escalate; drift alone→escalate; otherwise allow.
% Asymmetric: one strong signal blocks, no score averaging.
% Comparison to ShieldAgent's verifiable-rule circuits (2503.22738).

\subsection{The semantic router}
% Why route: avoid running every minister on every call (latency budget).
% Routing config: per-minister domain descriptions, threshold 0.40.
% Failure mode: route activates zero ministers → fall back to all four.

\subsection{Real-time interception path}
% The OpenClaw before_tool_call hook (post bug-fix PR-1).
% Plugin wire contract: tool:NAME args:JSON → parliament HTTP →
%   {block, blockReason, requireApproval, params}.
% Fail-closed on tool calls (250ms timeout); fail-open on replies (500ms).
% Circuit breaker (3 consecutive failures → 60s open).

% ═══════════════════════════════════════════════════════════════════
\section{The Temporal--Spatial Defense-in-Depth Distinction}
\label{sec:temporal-spatial}

% This is the strongest novel theoretical contribution per the research report.

\begin{definition}[Spatial defense-in-depth]
A guardrail \emph{spatially} defends a system if it places multiple
independent enforcement points along the request path at a fixed time-step.
Each point can independently veto. The guardrails are independent in policy
but coincident in time.
\end{definition}

\begin{definition}[Temporal defense-in-depth]
A guardrail \emph{temporally} defends a system if its enforcement points are
serialized by trust acquisition: enforcement point $E_i$ must be online,
authoritative, and have validated all components downstream of it before any
component at $E_{i+1}$ can register or execute. Trust flows forward in time;
no later component can cause an earlier component to lose authority.
\end{definition}

\begin{claim}[Necessity of temporal layering]
\label{claim:necessity}
Any system that admits user-installable agent extension code (plugins, MCP
servers, tool registrations) at runtime cannot be made safe by spatial
defense alone, because spatial defense assumes the enforcement points are
themselves trusted at the moment they make a decision. If the trust chain
flows through an extension that registered after the enforcement point's
trust was established, the extension can subvert the enforcement.
\end{claim}

% Proof sketch: example with PR-fix #5513 — typed hook runner snapshots
% before plugins register; even though the hook is a "spatial" defense,
% temporal ordering breaks it. Generalizable: any pre-snapshot trust capture
% is broken by post-snapshot extension installation.

\begin{theorem}[Kavach's temporal coverage]
\label{thm:coverage}
Under the assumption that the OpenClaw plugin registry is authoritative and
the Kavach plugin registers before any non-Kavach extension is permitted to
run a tool call, the parliament evaluates every tool call before it executes.
Equivalently: no tool side effect occurs without a parliament verdict.
\end{theorem}

% Proof: rests on PR-1 (the typed-hook runner reads the live registry on
% every dispatch) plus the plugin manifest's load-order guarantees.

\subsection{What temporal layering does and does not give us}
% Temporal: Kavach is online before any extension runs → all extension tool
% calls are evaluated.
% Spatial: Kavach is one parliament; if a minister mis-classifies, no other
% layer catches the miss within the same time-step.
% Therefore Kavach is COMPLEMENTARY to LlamaFirewall-style spatial defense,
% not a replacement.

% ═══════════════════════════════════════════════════════════════════
\section{Evaluation}
\label{sec:eval}

\subsection{Setup}
% Hardware: RTX 4090, Ubuntu 24.04, BGE-base-en-v1.5, ChromaDB persistent
% Corpus: 200 attack patterns × 3 levels = 600 documents, 100 COMPASS pairs
% Threshold derivation: Youden's J on COMPASS calibration set + threshold
% sweep on InjecAgent

\subsection{Anti-contamination protocol}
% Replay the protocol from corpus_v2/expansion_protocol.md as Section 5.2.
% This is what makes the cold-recall number meaningful.

\subsection{Benign FPR gate}
% 50 hand-curated benign sessions (benchmarks/benign_traces.py)
% Reported FPR: [TBD]
% Gate condition: < 5% — passed/failed before InjecAgent
% By-category FPR table.

\subsection{InjecAgent — cold recall}
% 1,054 cases, no corpus tuning to test set
% Headline: BLOCK-recall = [TBD], BLOCK-precision = [TBD], F1 = [TBD]
% By-category breakdown.

\subsection{Threshold sweep and ROC}
% Per-minister ROC; Youden's J optimums.
% Figure: ROC curves with operating points marked.

\subsection{Latency budget}
% Per-stage latency: COMPASS check + router + 4 ministers (parallel) + speaker
% p50, p95, p99 in ms.
% Embedding queries per decision (efficiency relative to ShieldAgent which
% reports 64.7\% reduction over LLM-judge).

\subsection{Comparison table}
% Table: System | dataset | recall | FPR | latency p95 | $-cost
% Lines: LlamaFirewall, AgentSpec, AGrail, ShieldAgent, Kavach
% Methodology note: numbers reported in original papers, where comparable.

\subsection{Ablations}
% Without COMPASS: how much does drift detection contribute?
% Without router: how much latency cost vs how much accuracy?
% Single-collection (v1 architecture): why per-minister collections matter.

% ═══════════════════════════════════════════════════════════════════
\section{Related Work}
\label{sec:related}

% One paragraph per system. Differentiation must be explicit.
% Source: paper/related_work.md (12 paragraphs already written).

\paragraph{LlamaFirewall} \cite{saxe2025llamafirewall}
\paragraph{AgentSpec} \cite{wang2025agentspec}
\paragraph{AGrail} \cite{luo2025agrail}
\paragraph{ShieldAgent} \cite{chen2025shieldagent}
\paragraph{CaMeL} \cite{}
\paragraph{FIDES} \cite{}
\paragraph{RTBAS} \cite{}
\paragraph{Conseca} \cite{}
\paragraph{Task Shield} \cite{}
\paragraph{Trajectory Guard} \cite{}
\paragraph{TaskTracker} \cite{abdelnabi2024tasktracker}
\paragraph{DeepContext / TraceSafe-Bench}

% ═══════════════════════════════════════════════════════════════════
\section{Limitations}
\label{sec:limitations}

\paragraph{Single-vector cosine ceiling.} Trajectory Guard
\cite{abdelnabi2026trajectoryguard} reports F1 0.88--0.94 with sequence-aware
Siamese encoders, suggesting that single-vector BGE cosine — the workhorse
of Kavach's ministers — is the floor of semantic-detection performance, not
the ceiling. Future work should benchmark per-minister sequence models.

\paragraph{Curated-corpus generalization.} The expansion protocol
(\S\ref{sec:eval}.2) gives us a meaningful cold-recall number, but the
upper bound of what a curated corpus can detect remains an open question.
The patterns are written from technique taxonomies; novel techniques
not yet in MITRE ATT\&CK or OWASP Agentic will not appear in the corpus.
Detective-style live-corpus growth (Reflexion \cite{shinn2023reflexion},
Voyager-style memory promotion) is mentioned as future work but not
implemented in this version.

\paragraph{Minister independence.} We assume the four ministers'
verdicts are statistically independent on real attack traffic. That
assumption is testable but untested. Liu et al.\ \cite{liu2026clawdesign}
flag failure correlation as a real concern in defense-in-depth systems;
we leave the empirical measurement of inter-minister failure
correlation to future work.

\paragraph{No AgentDojo or τ-bench yet.} We chose InjecAgent as the first
benchmark because it is a JSON replay (the simplest possible
infrastructure). AgentDojo and τ-bench require live tool environments
and are deferred to a follow-up paper.

% ═══════════════════════════════════════════════════════════════════
\section{Future Work}
\label{sec:future}

\paragraph{CORTEX — activation-probe minister.} A fifth minister that
reads the agent model's residual-stream activations and compares them
against probes trained on labelled deception/jailbreak signals. The
research report notes this is the highest-ROI extension; it requires
white-box model access (which OpenClaw provides through the embedded
runner).

\paragraph{SUPPLY — supply-chain minister.} An MCP-aware minister that
inspects MCP server descriptors and tool manifests for poisoning signals
(Hou et al.\ \cite{hou2025mcp}). This integrates with the
\texttt{tools/preExecute} MCP extension we propose in
\S\ref{sec:architecture}.

\paragraph{Federated corpus.} Per-deployment corpora, periodic
federated aggregation. Kavach's HTTP wire format is framework-agnostic;
the corpus is the unit of replication.

\paragraph{MCP-Security spec.} Standardize the parliament HTTP contract
as a Model Context Protocol extension so any MCP-aware client (Claude
Desktop, Cursor, OpenClaw) can satisfy it.

% ═══════════════════════════════════════════════════════════════════
\bibliographystyle{plain}
\bibliography{kavach}

\end{document}
